\documentclass[11pt,a4paper]{article}
\usepackage[utf8]{inputenc}
\usepackage[T1]{fontenc}
\IfFileExists{lmodern.sty}{\usepackage{lmodern}}{\IfFileExists{mathptmx.sty}{\usepackage{mathptmx}}{}}
\usepackage[ngerman]{babel}
\usepackage[a4paper,margin=27mm]{geometry}
\usepackage{amsmath,amssymb}
\usepackage{booktabs}
\usepackage{array}
\renewcommand{\floatpagefraction}{0.85}
\renewcommand{\textfraction}{0.1}
\IfFileExists{enumitem.sty}{\usepackage{enumitem}\setlist[itemize]{leftmargin=1.5em,itemsep=2pt}}{}
\usepackage{tikz}
\usetikzlibrary{arrows.meta,positioning,calc}
\IfFileExists{microtype.sty}{\usepackage{microtype}}{}
\usepackage[hidelinks]{hyperref}
\setlength{\parskip}{0.35em}
\setlength{\parindent}{1.2em}
\newcommand{\SO}{Subjekt-Objekt}
\newcommand{\Werke}[1]{(W~2, S.~#1)}
\newcolumntype{R}[1]{>{\raggedright\arraybackslash}p{#1}}

\title{Zeit und Ewigkeit der Rekursion\\[0.4em]
\large Der Fortgang der Selbstkonstruktion der Identität zur Totalität\\ in Hegels Differenzschrift (1801)}
\author{}
\date{}

\begin{document}
\maketitle

\section{Die offene Frage}

In der Vergleichung des Schellingschen Prinzips mit dem Fichteschen stehen zwei Sätze, die einander zu widersprechen scheinen. Der eine: Jeder Teil des Subjekts und jeder Teil des Objekts ist selbst im Absoluten, eine Identität des Subjekts und Objekts; die Teilung geht ins Unendliche, und jeder Staub ist eine Organisation \Werke{96~f.}. Das Ganze ist in jedem Teil schon da. Der andere: Die beiden Wissenschaften erscheinen als der Fortgang der Entwicklung oder Selbstkonstruktion der Identität zur Totalität \Werke{111}. Das Ganze wird erst konstruiert, es gibt einen Fortgang, eine Entwicklung, ein Vorher und Nachher. Wenn die Totalität in jedem Teil ist, wozu ein Fortgang zu ihr? Und wenn sie erst am Ende eines Fortgangs steht, wie ist sie in jedem Teil?

Die Frage ist die nach der Zeit. Ein Fortgang ist in der Zeit; ein Ganzes, das in jedem Teil ist, ist es nicht. Die Differenzschrift hat zu dieser Frage mehr zu sagen, als ihre Anlage vermuten lässt, denn die Zeit ist in ihr der Ort, an dem Fichtes System scheitert, und darum wird an mehreren Stellen bestimmt, was Zeit ist, was ihr Aufheben ist und was Ewigkeit heißt. Diese Bestimmungen werden im Folgenden gesammelt, dann an der Struktur der Rekursion nachgeprüft, die die beiden Sätze teilen, und schließlich in die Natur und in den Geist verfolgt.

Zitiert wird nach Band~2 der Theorie-Werkausgabe (Suhrkamp 1986), abgekürzt W~2, mit der dort eingedruckten Paginierung.

%=====================================================================
\section{Hegels Bestimmungen der Zeit}

\subsection{Der Zeit-Progress}

Die erste Bestimmung steht in der Vorerinnerung, als Diagnose. Fichtes Prinzip, die transzendentale Anschauung, Ich~=~Ich, verliert sich, sobald das System sich bildet, in einen \glqq Zeit-Progreß in infinitum\grqq{}, und das reine Bewusstsein erhält dadurch die Form eines \glqq durch eine objektive Unendlichkeit bedingten\grqq{} \Werke{11}. Die Zeit tritt hier als das auf, worin die Identität sich verliert: als die Reihe, die nicht endet, und an deren Ende die Identität stünde, wenn sie endete.

\subsection{Zufälligkeit in der Zeit}

Die zweite Bestimmung steht im Kapitel über das Bedürfnis der Philosophie und ist allgemeiner. Das Absolute muss sich als objektive Totalität setzen, und diese Objektivität hat eine zufällige Seite: \glqq Die Zufälligkeit ist eine Zufälligkeit in der Zeit, insofern die Objektivität des Absoluten als ein Fortgehen in der Zeit angeschaut wird; insofern sie aber als Nebeneinander im Raum erscheint, ist die Entzweiung klimatisch\grqq{} \Werke{22}. Zeit und Raum sind hier die beiden Formen, in denen das Absolute, sofern es objektiv wird, zufällig wird. Das Fortgehen in der Zeit ist nicht die Struktur des Absoluten, sondern die Form, in der seine Objektivität angeschaut wird, und in dieser Form ist sie kontingent: Was in der Zeit nacheinander kommt, könnte auch anders nacheinander kommen.

\subsection{Der Progress als Vermischung}

Die dritte Bestimmung steht im Kapitel über die transzendentale Anschauung und ist die genaueste. Der unendliche Progress, den Fichte und Kant als Idee postulieren, ist \glqq eine Vermischung von Empirischem und Vernünftigem; jenes ist die Anschauung der Zeit, dies die Aufhebung aller Zeit, die Verunendlichung derselben; im empirischen Progreß ist sie aber nicht rein verunendlicht, denn sie soll in ihm als Endliches, als beschränkte Momente, bestehen, -- er ist eine empirische Unendlichkeit\grqq{} \Werke{44}. Der Progress mischt also zwei Dinge: die Zeit, die aus beschränkten Momenten besteht, und ihre Aufhebung, die Unendlichkeit. Er hält beides nebeneinander, den Moment und das Unendliche, und darum ist er weder das eine noch das andere.

Dann folgt der entscheidende Satz. \glqq Die wahre Antinomie, die beides, das Beschränkte und Unbeschränkte, nicht nebeneinander, sondern zugleich als identisch setzt, muß damit zugleich die Entgegensetzung aufheben; indem die Antinomie die bestimmte Anschauung der Zeit postuliert, muß diese -- beschränkter Moment der Gegenwart und Unbeschränktheit seines Außersichgesetztseins -- beides zugleich, also Ewigkeit sein\grqq{} \Werke{44}. Ewigkeit ist hier definiert, und die Definition ist nicht die Negation der Zeit. Ewigkeit ist der beschränkte Moment der Gegenwart und die Unbeschränktheit zugleich. Nicht nacheinander, wie im Progress, wo auf jeden Moment ein nächster folgt und die Unbeschränktheit außerhalb liegt; sondern in einem: die Gegenwart, die ihr Außersich in sich hat.

\subsection{Das wahre Aufheben der Zeit}

Die vierte Bestimmung steht in der Darstellung des Fichteschen Systems, und sie führt die dritte aus. Fichtes Streben ist eine Forderung der Vereinigung, \glqq die aber nicht geschehen soll\grqq{}, und drückt sich im unendlichen Progress aus \Werke{70}. \glqq Das in die Ewigkeit verlängerte Dasein schließt beides, Unendlichkeit der Idee und Anschauung in sich, aber beides in solchen Formen, die ihre Synthese unmöglich macht\grqq{} \Werke{70}. Die Zeit \glqq schließt unmittelbar Entgegensetzung, ein Außereinander in sich, und das Dasein in der Zeit ist ein sich Entgegengesetztes, Mannigfaltiges, und die Unendlichkeit ist außer ihr\grqq{} \Werke{70}. Dann eine Bemerkung, die hier nur festgehalten wird: Der Raum sei gleichfalls ein Außersichgesetztsein, könne aber \glqq in seinem Charakter der Entgegensetzung eine unendlich reichere Synthese genannt werden als die Zeit\grqq{} \Werke{70}. Der Vorzug, den die Zeit für den Progress erhält, liege nur darin, dass das Streben als ein Inneres gesetzt wird und Ich als Punkt, als Seele hypostasiert.

Und dann: \glqq Soll die Zeit Totalität sein, als unendliche Zeit, so ist die Zeit selbst aufgehoben, und es war nicht nötig, zu ihrem Namen und zu einem Progreß des verlängerten Daseins zu fliehen. Das wahre Aufheben der Zeit ist zeitlose Gegenwart, d.\,i. Ewigkeit; und in dieser fällt das Streben und das Bestehen absoluter Entgegensetzung weg\grqq{} \Werke{71}. Zwei Ewigkeiten stehen einander gegenüber: das in die Ewigkeit verlängerte Dasein, das nur die Zeit ohne Ende ist, und die zeitlose Gegenwart, die die Zeit aufhebt. Die erste ist der Progress, die zweite die Anschauung.

Zwei Seiten später wird die Konsequenz für das Ich gezogen: Es \glqq produziert in dem unendlichen Progreß des verlängerten Daseins endlos Teile von sich, aber nicht sich selbst in der Ewigkeit des Sich-selbst-Anschauens als \SO{}\grqq{} \Werke{72}. Und noch einmal, am Ende der Darstellung des praktischen Teils: Wo die höchste Synthese erwartet wird, bleibt \glqq immer dieselbe Antithese der beschränkten Gegenwart und einer außer ihr liegenden Unendlichkeit\grqq{}; Ich~=~Ich sei das Absolute, \glqq aber so weit bringt es Ich im System nicht und, wenn die Zeit eingemischt werden soll, nie\grqq{} \Werke{76}.

\subsection{Der Fortgang}

Die fünfte Bestimmung ist die, von der die Frage ausging. Die beiden Wissenschaften streben nach dem Indifferenzpunkt; er liegt als Identität überall in ihnen, als absolute Totalität außer ihnen; sie sind die Linien, die den Pol mit dem Mittelpunkt verknüpfen; der Mittelpunkt ist gedoppelt, Identität und Totalität, \glqq und insofern erscheinen beide Wissenschaften als der Fortgang der Entwicklung oder Selbstkonstruktion der Identität zur Totalität\grqq{} \Werke{111}. Der Fortgang ist der der Wissenschaften, nicht des Absoluten; er ist der Weg von der Identität, die in jedem Teil liegt, zur Totalität, die das Ganze ist; und die Wissenschaften sind dieser Weg, die Linien.

Die Bestimmungen lassen sich in einer Reihe ordnen: Zeit ist Außereinander, Mannigfaltigkeit beschränkter Momente, mit der Unendlichkeit außer sich. Der Progress ist die Vermischung von Zeit und Unendlichkeit, die beide nebeneinander hält. Die Ewigkeit ist die Gegenwart, die ihre Unbeschränktheit in sich hat, zeitlose Gegenwart. Und der Fortgang ist der Weg der Wissenschaft von der Identität zur Totalität. Was fehlt, ist die Angabe, wie sich Ewigkeit und Fortgang zueinander verhalten. Dazu die Struktur.

%=====================================================================
\section{Dieselbe Rekursion in zwei Weisen}

\subsection{Die Stufen}

Die Rekursion der Differenzschrift ist ein Baum. Das Absolute ist \SO{}; seine beiden Seiten sind je \SO{}; jede Seite jeder Seite wieder; ohne Ende. Man kann diesen Baum auf zwei Weisen haben, und die beiden Weisen sind die beiden Ewigkeiten Hegels.

Die erste Weise baut ihn auf. Man beginnt mit der Wurzel, setzt ihre beiden Seiten, setzt deren Seiten, und so fort, Stufe um Stufe. Auf jeder Stufe hat man einen endlichen Baum. Diese Bäume lassen sich nachrechnen. Auf der Stufe $n$ hat der Baum $2^n$ Blätter, und jedes Blatt ist ein Rand: ein Knoten ohne Seiten, ein \SO{}, das noch nicht in Subjekt und Objekt geteilt ist, also ein Pol, kein Ganzes. Der Anteil solcher Randknoten an allen Knoten des Baumes ist auf der ersten Stufe zwei Drittel und nähert sich mit wachsender Stufe der Hälfte, ohne sie je zu unterschreiten. Auf keiner Stufe ist der Baum ohne Rand; auf jeder Stufe ist mindestens die Hälfte seiner Knoten Rand. Und das Ganze, der Baum ohne Rand, ist auf keiner Stufe erreicht, denn jede Stufe ist endlich.

Das ist der Zeit-Progress. Jede Stufe ist ein beschränkter Moment; die nächste Stufe liegt außer ihr; das Ganze liegt außer allen. Das Ich, das den Baum so aufbaut, \glqq produziert in dem unendlichen Progreß des verlängerten Daseins endlos Teile von sich, aber nicht sich selbst\grqq{}. Die Stufen sind die Teile; das Selbst wäre der Baum ohne Rand; und der ist in keiner Stufe. Man kann auch sagen, warum er in keiner Stufe sein kann: Wäre der Baum eine Stufe, also endlich, so hätte er eine Größe $s$, und diese müsste, weil er aus einer Wurzel und zwei Bäumen derselben Größe besteht, die Gleichung $s = 1 + 2s$ erfüllen. Die einzige Lösung ist $s = -1$. Das Ganze der Rekursion ist keine Anzahl. Das ist die Rechnung zu Hegels Satz, das Absolute sei \glqq kein Quantum, sondern Totalität\grqq{} \Werke{99}.

\subsection{Der Selbstbezug}

Die zweite Weise baut ihn nicht auf, sondern gibt ihn durch das, was er ist. Der Baum $T$ ist das, dessen beide Seiten $T$ sind: $T = (T, T)$. Das ist keine Anweisung zum Aufbauen, sondern eine Bestimmung durch Selbstbezug. Sie sagt nicht, wie der Baum entsteht, sondern was an jeder Stelle des Baumes gilt, und sie sagt es für alle Stellen auf einmal. Ein so bestimmter Baum hat keinen Rand: Jeder Knoten hat beide Seiten, weil jeder Knoten $T$ ist. Und jeder Teilbaum ist dem Ganzen gleich; man kann an einem beliebigen Pfad beliebig weit hinabsteigen, der Teilbaum, den man dort findet, entfaltet sich genau wie die Wurzel. Das ist geprüft worden, an zweihundert zufälligen Pfaden bis zur Tiefe vierzig, und es gilt nicht ungefähr, sondern exakt, weil der Teilbaum nicht nur gleich der Wurzel, sondern die Wurzel ist.

Das ist die zeitlose Gegenwart. An jeder Stelle ist das Ganze da, nicht als Ergebnis eines Aufstiegs, sondern weil die Stelle durch das Ganze bestimmt ist. Es gibt keinen Rand, also keine Stelle, die noch nicht geteilt wäre, also kein Sollen. \glqq In dieser fällt das Streben und das Bestehen absoluter Entgegensetzung weg\grqq{} \Werke{71}. Und es gibt kein Nacheinander der Stufen, weil es keine Stufen gibt; die Teilung ist nicht ein Vorgang, sondern eine Eigenschaft jedes Knotens.

Der Unterschied der beiden Weisen ist nicht ein Unterschied dessen, was da ist. Jede endliche Stufe ist ein Anfangsstück des selbstbezüglichen Baumes, und der selbstbezügliche Baum ist das, worauf die Stufen zulaufen. Der Unterschied ist der der Gegebenheit: einmal durch Aufbau in der Folge, einmal durch Bestimmung im Selbstbezug. Hegel hat für die beiden Weisen die Worte Reflexion und Anschauung. Die Reflexion setzt Stufe um Stufe, in Sätzen, und findet an jedem Satz den nächsten; die Anschauung hat das Ganze, in dem alle Sätze schon gelten. Was der Progress als eine unendliche Aufgabe vor sich hat, hat die Anschauung als Bestimmung in sich.

\begin{figure}[!htb]
\centering
\resizebox{\linewidth}{!}{%
\begin{tikzpicture}[
  every node/.style={align=center, font=\small},
  nd/.style={draw, circle, inner sep=1.6pt, fill=white},
  rd/.style={draw, circle, inner sep=1.6pt, fill=gray!50},
  lbl/.style={font=\scriptsize}]
% --- links: Stufen
\node[font=\small\bfseries] at (-4.2,2.6) {Aufbau in Stufen (Progress)};
\foreach \y/\n in {1.6/1, 0.3/2, -1.2/3} {
  \node[lbl, anchor=east] at (-7.4,\y) {Stufe \n};
}
% Stufe 1
\node[nd] (a0) at (-4.2,1.9) {}; \node[rd] (a1) at (-4.9,1.3) {}; \node[rd] (a2) at (-3.5,1.3) {};
\draw (a0)--(a1); \draw (a0)--(a2);
% Stufe 2
\node[nd] (b0) at (-4.2,0.7) {}; \node[nd] (b1) at (-5.0,0.2) {}; \node[nd] (b2) at (-3.4,0.2) {};
\node[rd] (b11) at (-5.4,-0.3) {}; \node[rd] (b12) at (-4.6,-0.3) {}; \node[rd] (b21) at (-3.8,-0.3) {}; \node[rd] (b22) at (-3.0,-0.3) {};
\draw (b0)--(b1); \draw (b0)--(b2); \draw (b1)--(b11); \draw (b1)--(b12); \draw (b2)--(b21); \draw (b2)--(b22);
% Stufe 3
\node[nd] (c0) at (-4.2,-0.8) {}; \node[nd] (c1) at (-5.1,-1.25) {}; \node[nd] (c2) at (-3.3,-1.25) {};
\node[nd] (c11) at (-5.55,-1.7) {}; \node[nd] (c12) at (-4.65,-1.7) {}; \node[nd] (c21) at (-3.75,-1.7) {}; \node[nd] (c22) at (-2.85,-1.7) {};
\foreach \p/\x in {c11/-5.55, c12/-4.65, c21/-3.75, c22/-2.85} {
  \node[rd] (\p a) at ({\x-0.22},-2.15) {}; \node[rd] (\p b) at ({\x+0.22},-2.15) {};
  \draw (\p)--(\p a); \draw (\p)--(\p b);
}
\draw (c0)--(c1); \draw (c0)--(c2); \draw (c1)--(c11); \draw (c1)--(c12); \draw (c2)--(c21); \draw (c2)--(c22);
\node[lbl] at (-4.2,-2.65) {$\vdots$};
\node[lbl, text width=5.2cm] at (-4.2,-3.4) {graue Knoten: Rand, noch ungeteilt.\\ Auf jeder Stufe mindestens die Hälfte der Knoten. Das Ganze in keiner Stufe: $s = 1 + 2s$ hat keine Anzahl als Lösung.};
% --- rechts: Selbstbezug
\node[font=\small\bfseries] at (3.9,2.6) {Bestimmung durch Selbstbezug (Anschauung)};
\node[draw, rounded corners=3pt, inner sep=6pt, fill=gray!12] (T) at (3.9,1.2) {$T$\\ \footnotesize \SO{}};
\node[draw, rounded corners=3pt, inner sep=5pt] (TL) at (2.3,-0.6) {$T$\\ \footnotesize Seite: Subjekt};
\node[draw, rounded corners=3pt, inner sep=5pt] (TR) at (5.5,-0.6) {$T$\\ \footnotesize Seite: Objekt};
\draw[-{Latex[length=2mm]}, thick] (T) -- (TL); \draw[-{Latex[length=2mm]}, thick] (T) -- (TR);
\draw[-{Latex[length=2mm]}, thick, dashed, gray!60!black] (TL.west) to[bend left=50] node[left, lbl] {ist $T$} (T.west);
\draw[-{Latex[length=2mm]}, thick, dashed, gray!60!black] (TR.east) to[bend right=50] node[right, lbl] {ist $T$} (T.east);
\node[lbl] at (3.9,-1.7) {$T = (T,\,T)$};
\node[lbl, text width=5.2cm] at (3.9,-3.4) {kein Rand, kein Nacheinander. Jeder Teilbaum ist nicht nur gleich der Wurzel, er ist die Wurzel. \glqq Jeder Staub eine Organisation.\grqq{}};
\end{tikzpicture}}
\caption{Dieselbe Rekursion in zwei Weisen. Links der Aufbau in Stufen: jede Stufe endlich, mit Rand, das Ganze außer allen Stufen (W~2, S.~11, 72). Rechts die Bestimmung durch Selbstbezug: das Ganze an jeder Stelle, zeitlose Gegenwart (W~2, S.~44, 71).}
\label{fig:zweiweisen}
\end{figure}

\subsection{Fortgang und Totalität}

Nun lässt sich die offene Frage beantworten. Der Fortgang der Selbstkonstruktion der Identität zur Totalität ist nicht der Aufbau des Baumes in Stufen. Wäre er das, so wäre er der Progress, und die Totalität läge außer ihm, und Hegel hätte gegen Schelling denselben Einwand wie gegen Fichte. Der Fortgang ist etwas anderes: Er ist der Weg durch den Baum, der schon ganz ist. Die Wissenschaft geht die Linie vom Pol zum Mittelpunkt, sie durchläuft die Knoten, sie stellt an jedem Knoten dar, dass er \SO{} ist. Der Baum wird dabei nicht erzeugt, sondern entfaltet; das Wort Entwicklung bei Hegel heißt genau dies, Auswicklung dessen, was eingewickelt schon da ist. Die Totalität ist die Struktur; der Fortgang ist ihr Durchgang.

Damit klärt sich, was die Zeit in der Differenzschrift ist. Zeit ist die Form des Durchgangs. Sie ist nicht die Form des Absoluten; das Absolute ist zeitlose Gegenwart. Aber die Wissenschaft, die das Absolute konstruiert, ist Reflexion, setzt Satz um Satz, und Satz um Satz ist Nacheinander. Die Zeit ist die Ordnung, in der die Knoten des Baumes besucht werden. Und diese Ordnung ist zufällig, in dem genauen Sinn, den Hegel angibt: \glqq Die Zufälligkeit ist eine Zufälligkeit in der Zeit, insofern die Objektivität des Absoluten als ein Fortgehen in der Zeit angeschaut wird\grqq{} \Werke{22}. Der Baum legt fest, was an jedem Knoten gilt; er legt nicht fest, in welcher Reihenfolge die Knoten durchlaufen werden. Man kann zuerst die Seite des Subjekts hinabsteigen oder zuerst die des Objekts; man kann in die Tiefe gehen oder in die Breite. Die Struktur ist gegen die Reihenfolge gleichgültig, und darum ist die Reihenfolge, die Zeit, das Zufällige an der Objektivität des Absoluten.

Das ist auch der Grund, warum es zwei Wissenschaften gibt und nicht eine. Die Transzendentalphilosophie und die Naturphilosophie sind zwei Durchgänge durch denselben Baum, der eine von der Seite des Subjekts her, der andere von der Seite des Objekts her. Beide durchlaufen alle Knoten; beide finden an jedem Knoten \SO{}; aber sie durchlaufen sie in entgegengesetzter Ordnung, und darum erscheinen sie einander entgegengesetzt, obwohl \glqq in beiden ein und ebendasselbe in den notwendigen Formen seiner Existenz konstruiert wird\grqq{} \Werke{101}. Ihr Widerspruch ist ein Widerspruch der Reihenfolge, nicht der Sache. Der höhere Standpunkt, \glqq der in beiden ebendasselbe Absolute erkennt\grqq{}, ist der, der den Baum hat und nicht nur einen Durchgang.

\begin{figure}[!htb]
\centering
\resizebox{0.92\linewidth}{!}{%
\begin{tikzpicture}[
  every node/.style={align=center, font=\small},
  nd/.style={draw, circle, inner sep=2pt, fill=white},
  ed/.style={thick},
  pathA/.style={-{Latex[length=1.8mm]}, very thick, gray!55!black},
  pathB/.style={-{Latex[length=1.8mm]}, very thick, dashed, gray!55!black},
  lbl/.style={font=\scriptsize}]
\node[nd, fill=gray!20] (r) at (0,2.6) {};
\node[nd] (s) at (-3.8,1.3) {}; \node[nd] (o) at (3.8,1.3) {};
\node[nd] (ss) at (-5.7,0) {}; \node[nd] (so) at (-1.9,0) {}; \node[nd] (os) at (1.9,0) {}; \node[nd] (oo) at (5.7,0) {};
\foreach \p/\x in {ss/-5.7, so/-1.9, os/1.9, oo/5.7} {
  \node[nd] (\p a) at ({\x-0.7},-1.1) {}; \node[nd] (\p b) at ({\x+0.7},-1.1) {};
  \draw[ed] (\p)--(\p a); \draw[ed] (\p)--(\p b);
  \draw[dotted, thick] (\p a) -- ++(0,-0.6); \draw[dotted, thick] (\p b) -- ++(0,-0.6);
}
\draw[ed] (r)--(s); \draw[ed] (r)--(o); \draw[ed] (s)--(ss); \draw[ed] (s)--(so); \draw[ed] (o)--(os); \draw[ed] (o)--(oo);
\node[lbl, above=3pt of r] {Absolutes: \SO{}, zeitlose Gegenwart};
\node[lbl, below left=1pt and -2pt of s] {Subjekt}; \node[lbl, below right=1pt and -2pt of o] {Objekt};
% Durchgang A: Transzendentalphilosophie (vom Subjektpol her), oberhalb der Kanten
\draw[pathA] ($(ss)+(-0.3,0.3)$) .. controls ($(s)+(-1.2,0.9)$) and ($(r)+(-1.6,0.6)$) .. ($(r)+(-0.35,0.25)$);
\draw[pathA] ($(r)+(0.35,0.25)$) .. controls ($(r)+(1.6,0.6)$) and ($(o)+(1.2,0.9)$) .. ($(oo)+(0.3,0.3)$);
% Durchgang B: Naturphilosophie (vom Objektpol her), unterhalb der Kanten
\draw[pathB] ($(oo)+(0.15,-0.4)$) .. controls ($(o)+(0.9,-0.9)$) and ($(r)+(1.2,-0.9)$) .. ($(r)+(0.2,-0.4)$);
\draw[pathB] ($(r)+(-0.2,-0.4)$) .. controls ($(r)+(-1.2,-0.9)$) and ($(s)+(-0.9,-0.9)$) .. ($(ss)+(-0.15,-0.4)$);
\node[lbl, text width=3.9cm] at (-5.4,-2.7) {durchgezogen: Transzendentalphilosophie, Durchgang vom Subjektpol her};
\node[lbl, text width=3.9cm] at (5.4,-2.7) {gestrichelt: Naturphilosophie, Durchgang vom Objektpol her};
\node[lbl, text width=5.6cm] at (0,-2.7) {Der Baum ist derselbe. Die Reihenfolge des Durchgangs ist die Zeit: \glqq Zufälligkeit in der Zeit\grqq{} (S.~22). Der Fortgang ist Entfaltung, nicht Erzeugung (S.~111).};
\end{tikzpicture}}
\caption{Fortgang und Totalität. Die beiden Wissenschaften sind zwei Durchgänge durch denselben, schon ganzen Baum, in entgegengesetzter Reihenfolge (W~2, S.~101, 111). Die Zeit ist die Ordnung des Durchgangs, nicht die Struktur des Durchlaufenen.}
\label{fig:durchgang}
\end{figure}

%=====================================================================
\section{Natur: der Umlauf}

Hegels Definition der Ewigkeit, der beschränkte Moment der Gegenwart und die Unbeschränktheit seines Außersichgesetztseins zugleich, hat in der Natur eine Gestalt, die genau so gebaut ist, und die Naturphilosophie der Zeit kannte sie: den Umlauf.

Eine Bewegung auf einer Geraden hat die Struktur des Progresses. Jeder Ort ist ein beschränkter Moment; der nächste liegt außer ihm; die Gerade endet nicht; und die Unendlichkeit, das Ende der Geraden, ist außer jedem Ort. Die Folge der Orte, die sich einem Grenzwert nähert, ohne ihn zu erreichen, ist die Bahn des Strebens.

Eine Bewegung auf einem Kreis hat die andere Struktur. Auch hier ist jeder Ort ein beschränkter Moment. Aber die Bahn kehrt zurück; sie ist unbeschränkt, denn sie hat kein Ende, und doch liegt ihre Unbeschränktheit nicht außer dem Ort, sondern der Ort hat sie in sich: Von jedem Punkt des Kreises aus ist der ganze Kreis bestimmt, das Gesetz der Bewegung ist an jedem Ort das ganze Gesetz, und kein Ort ist dem Ganzen näher als ein anderer. Der Umlauf ist die Zeit, die in sich zurückgeht. Er ist Zeit, denn er hat Nacheinander; und er ist Ewigkeit, denn jeder seiner Momente enthält das Ganze. Der Moment der Gegenwart ist beschränkt, ein Punkt der Bahn; seine Unbeschränktheit, die Bahn, ist in ihm; beides zugleich.

Der Magnet, den Hegel an der Stelle über die Pole der Indifferenz selbst nennt \Werke{110}, zeigt, wie aus dem Umlauf Ruhe wird. Ein Kreisstrom ist eine Bewegung: Ladung läuft um, in der Zeit, an jedem Moment an einem anderen Ort. Das Feld, das dieser Umlauf erzeugt, ist ruhend. Es ändert sich nicht, obwohl das, was es erzeugt, sich ununterbrochen ändert. Die beiden Pole und die Indifferenz zwischen ihnen sind da, unbewegt, als das Ganze des Umlaufs; sie sind die zeitlose Gegenwart der bewegten Ladung. Und wer den Magneten teilt, findet in jedem Stück wieder beide Pole und die Indifferenz, weil in jedem Stück der Umlauf ganz ist. Der Umlauf gebiert die Pole; die Pole sind, woran man den Umlauf erkennt.

Das ist die Gestalt, die Hegels Satz in der Natur hat: Das Absolute ist die zeitlose Gegenwart eines Umlaufs, der die Zeit als sein Moment in sich hat. Die Zeit ist nicht abwesend; sie ist aufgehoben, im doppelten Sinn, den das Wort bei Hegel trägt, beseitigt und bewahrt. Was im Progress als unendliche Gerade auseinandergezogen ist, ist im Umlauf in sich zurückgebogen, und in dieser Rückbiegung wird die Zeit zur Ewigkeit, ohne aufzuhören, Bewegung zu sein.

Ein Vorbehalt ist nötig. Der Umlauf ist ein Gegenstand der Naturphilosophie, also der Wissenschaft vom objektiven \SO{}; er ist ein Beispiel, nicht das Absolute. Aber er ist das Beispiel, das die Definition der Ewigkeit auf S.~44 wörtlich erfüllt, und er ist das Beispiel, an dem sich zeigt, dass Ewigkeit nicht Stillstand heißt.

\begin{figure}[!htb]
\centering
\resizebox{0.95\linewidth}{!}{%
\begin{tikzpicture}[every node/.style={align=center, font=\small}, lbl/.style={font=\scriptsize}]
% links: Gerade
\node[font=\small\bfseries] at (-4.6,2.4) {Gerade: Progress};
\draw[thick, -{Latex[length=2mm]}] (-7.6,0.6) -- (-1.6,0.6);
\foreach \i in {1,...,9} {
  \pgfmathsetmacro{\x}{-7.6 + 5.4*(1 - 1/\i)}
  \fill (\x,0.6) circle (1.8pt);
}
\draw[fill=white, thick] (-2.2,0.6) circle (2.4pt);
\node[lbl] at (-2.2,1.1) {Grenze:\\ außer jedem Ort};
\node[lbl] at (-7.6,0.15) {Anfang};
\node[lbl, text width=5.6cm] at (-4.6,-0.9) {jeder Ort ein beschränkter Moment; die Unendlichkeit außer ihm; \glqq endlos Teile von sich, aber nicht sich selbst\grqq{} (S.~72)};
% rechts: Kreis
\node[font=\small\bfseries] at (4.6,2.4) {Kreis: Umlauf, zeitlose Gegenwart};
\draw[thick] (4.6,0.3) circle (1.5cm);
\draw[-{Latex[length=2mm]}, thick] (4.6,1.8) arc (90:60:1.5cm);
\foreach \a in {0,45,...,315} { \fill ({4.6+1.5*cos(\a)},{0.3+1.5*sin(\a)}) circle (1.6pt); }
\fill[gray!60!black] ({4.6+1.5*cos(20)},{0.3+1.5*sin(20)}) circle (2.6pt);
\draw[gray] ({4.6+1.5*cos(20)+0.1},{0.3+1.5*sin(20)+0.05}) -- (6.7,1.5);
\node[lbl, anchor=west] at (6.75,1.5) {Gegenwart:\\ beschränkter Moment};
\node[lbl] at (4.6,0.3) {die Bahn:\\ an jedem Ort\\ ganz bestimmt};
\node[lbl, text width=5.6cm] at (4.6,-1.85) {\glqq beschränkter Moment der Gegenwart und Unbeschränktheit seines Außersichgesetztseins, beides zugleich, also Ewigkeit\grqq{} (S.~44)};
\end{tikzpicture}}
\caption{Gerade und Kreis als die beiden Gestalten der Zeit. Links das in die Ewigkeit verlängerte Dasein: die Grenze liegt außer jedem Ort. Rechts der Umlauf: die ganze Bahn ist in jedem Moment bestimmt; die Zeit ist aufgehoben und bewahrt.}
\label{fig:kreis}
\end{figure}

%=====================================================================
\section{Geist: die Ewigkeit des Sich-selbst-Anschauens}

Der Satz, an dem Hegel Fichtes Scheitern festmacht, nennt das, was Fichte verfehlt, mit einem Namen: die \glqq Ewigkeit des Sich-selbst-Anschauens als \SO{}\grqq{} \Werke{72}. Das ist die zeitlose Gegenwart, in den Raum des Geistes übersetzt. Und die Übersetzung lässt sich an der Struktur nachvollziehen.

Ich~=~Ich ist Selbstbezug. Das Ich ist das, was sich selbst setzt, dessen Gesetztes es selbst ist: Ich~=~(Ich, Ich), das Subjekt und sein Objekt, und beide dasselbe. Das ist die Bestimmung durch Selbstbezug, dieselbe Gestalt wie $T = (T, T)$. Fichte hat diese Bestimmung an den Anfang gestellt, und Hegel bestreitet nicht, dass sie das Absolute ist; er sagt ausdrücklich, das Selbstbewusstsein sei \glqq nicht erst auf das Absolute als ein Höheres bezogen, sondern es ist selbst das Absolute\grqq{} \Werke{96}. Was Hegel bestreitet, ist die Weise, in der Fichte den Selbstbezug dann hat. Fichte nimmt ihn nicht als Bestimmung, die an jeder Stelle gilt, sondern als Anweisung, die Stufe um Stufe ausgeführt wird: Das Ich setzt sich; es setzt ein Nicht-Ich; es setzt beide teilbar; es strebt, das Nicht-Ich zu bestimmen; und so fort. Die Rekursion wird als Aufbau in Stufen genommen, jede Stufe ein Akt, jeder Akt in der Zeit, und das Ganze, das Ich~=~Ich, das am Anfang stand, wird zum Ende, das nie erreicht wird. Der Selbstbezug, als Folge ausgeführt, ist der Progress; als Bestimmung angeschaut, ist er die Ewigkeit. Fichte hat den Selbstbezug richtig gesetzt und falsch gehabt.

Was heißt es dann, das Sich-selbst-Anschauen in der Ewigkeit zu haben? Es heißt, dass jede Reflexion, jedes Wissen von einem Wissen, jede Stufe des Sich-auf-sich-Richtens, nicht eine Annäherung an die Identität ist, sondern die Identität selbst, ganz. \glqq Jedes Erkennen eine Wahrheit\grqq{} \Werke{97}. Der Geist, der sich reflektiert, produziert nicht Teile von sich, sondern ist an jedem Ort der Reflexion er selbst. Die Reflexion ist Zeit, denn sie ist Akt, sie hat ein Vorher und ein Nachher. Aber sie ist Umlauf, nicht Gerade: Sie geht vom Ich zum Ich, sie kehrt zurück, und das, wozu sie zurückkehrt, ist nicht ein Späteres, sondern dasselbe. Die Gegenwart des Selbstbewusstseins ist ein beschränkter Moment, dieser Akt, jetzt; und sie hat die Unbeschränktheit in sich, denn dieser Akt ist das ganze Ich. Beides zugleich, also Ewigkeit.

Und damit ist auch die Zeit im Geist an ihrem Ort. Sie ist die Form des Durchgangs: die Ordnung, in der ein endliches Bewusstsein die Knoten seiner selbst besucht, was es zuerst weiß und was danach, welche Seite es zuerst reflektiert. Diese Ordnung ist zufällig, wie Hegel es von der Objektivität des Absoluten sagt; sie ist Bildung, Geschichte, Biographie. Was nicht zufällig ist, ist das, was an jedem Knoten gilt. Die zeitlose Gegenwart des Geistes ist nicht ein Ort außerhalb seiner Geschichte; sie ist das, was in jedem Moment seiner Geschichte ganz da ist, wenn er sich in ihm anschaut statt sich in ihm zu suchen.

%=====================================================================
\section{Ergebnis}

Die offene Frage war, wie die Totalität in jedem Teil sein kann, wenn es einen Fortgang zu ihr gibt. Die Antwort der Differenzschrift, aus ihren Bestimmungen der Zeit zusammengesetzt, lautet so.

Die Rekursion des \SO{} kann in zwei Weisen gehabt werden: als Aufbau in Stufen und als Bestimmung durch Selbstbezug. Der Aufbau ist der Progress: jede Stufe endlich, mit Rand, das Ganze außer allen Stufen, und das Ganze keine Anzahl. Der Selbstbezug ist die zeitlose Gegenwart: das Ganze an jeder Stelle, kein Rand, kein Sollen. Hegel nennt die erste Weise das in die Ewigkeit verlängerte Dasein und die zweite Ewigkeit; die erste ist Reflexion, die zweite Anschauung.

Der Fortgang der Selbstkonstruktion ist nicht der Aufbau, sondern der Durchgang: die Wissenschaft durchläuft den Baum, der schon ganz ist, und stellt an jedem Knoten dar, was an ihm gilt. Die Zeit ist die Ordnung dieses Durchgangs, und sie ist zufällig, weil die Struktur gegen die Reihenfolge gleichgültig ist. Die beiden Wissenschaften sind zwei Durchgänge in entgegengesetzter Ordnung durch denselben Baum.

Ewigkeit heißt bei Hegel nicht Abwesenheit der Zeit, sondern die Gegenwart, die ihre Unbeschränktheit in sich hat. In der Natur hat das die Gestalt des Umlaufs, der Zeit ist und in jedem Moment ganz; der Magnet ist die Ruhe eines solchen Umlaufs. Im Geist hat es die Gestalt des Sich-selbst-Anschauens, in dem jede Reflexion das ganze Ich ist. Fichte hat den Selbstbezug als Folge ausgeführt und dadurch in den Progress verwandelt; Hegel hält ihn als Bestimmung fest und findet ihn dadurch in jedem Moment.

Was offen bleibt, ist der Raum. Hegel nennt ihn ein Außersichgesetztsein wie die Zeit, aber \glqq eine unendlich reichere Synthese\grqq{} als sie \Werke{70}. Ob das Nebeneinander, das die andere Form der Zufälligkeit des Absoluten ist \Werke{22}, zur Rekursion in einem anderen Verhältnis steht als das Nacheinander, wird hier nicht entschieden.

\clearpage
\section*{Apparat}

\subsection*{1. Statustafel}
\begin{center}
\renewcommand{\arraystretch}{1.3}
\begin{tabular}{@{}R{8.6cm}R{2.6cm}R{3.4cm}@{}}
\toprule
Aussage & Status & Beleg \\
\midrule
Ewigkeit ist definiert als beschränkter Moment der Gegenwart und Unbeschränktheit zugleich; das wahre Aufheben der Zeit ist zeitlose Gegenwart. & belegt & W~2, S.~44, 71 \\
Zwei Ewigkeiten: das in die Ewigkeit verlängerte Dasein (Progress) und die zeitlose Gegenwart (Anschauung). & belegt & W~2, S.~70~f. \\
Zeit ist die Form der Zufälligkeit der Objektivität des Absoluten. & belegt & W~2, S.~22 \\
Der Fortgang ist der der Wissenschaften, von der Identität zur Totalität. & belegt & W~2, S.~111 \\
Die endlichen Stufen enthalten stets Rand (mindestens die Hälfte der Knoten); das Ganze ist keine Anzahl ($s = 1 + 2s$ ohne positive Lösung); der selbstbezügliche Baum hat jeden Teilbaum als Wurzel. & rechnerisch geprüft & Abschnitt 3 \\
Aufbau in Stufen $=$ Progress; Bestimmung durch Selbstbezug $=$ zeitlose Gegenwart. & Deutung, gestützt & W~2, S.~44, 71, 72 \\
Der Fortgang ist Durchgang durch den ganzen Baum, nicht Aufbau; die Zeit ist die Ordnung des Durchgangs. & Deutung & W~2, S.~22, 101, 111 \\
Die beiden Wissenschaften als Durchgänge in entgegengesetzter Ordnung. & Deutung, gestützt & W~2, S.~101 \\
Der Umlauf erfüllt die Definition der Ewigkeit; das Magnetfeld als Ruhe eines Umlaufs. & Deutung; physikalisch korrekt für den stationären Kreisstrom & W~2, S.~44, 110 \\
Fichte hat den Selbstbezug als Folge ausgeführt, Hegel als Bestimmung. & Deutung, gestützt & W~2, S.~72, 96 \\
Verhältnis des Raums zur Rekursion. & offen & W~2, S.~22, 70 \\
\bottomrule
\end{tabular}
\end{center}

\subsection*{2. Abbildungen}
\begin{itemize}
\item Abbildung~\ref{fig:zweiweisen}: Dieselbe Rekursion als Aufbau in Stufen und als Bestimmung durch Selbstbezug.
\item Abbildung~\ref{fig:durchgang}: Die beiden Wissenschaften als Durchgänge durch denselben Baum; die Zeit als Ordnung des Durchgangs.
\item Abbildung~\ref{fig:kreis}: Gerade und Kreis als Gestalten des Progresses und der zeitlosen Gegenwart.
\end{itemize}

\subsection*{3. Sachen}
\begin{itemize}
\item Zeit als Außereinander beschränkter Momente
\item Zeit-Progress in infinitum
\item in die Ewigkeit verlängertes Dasein
\item zeitlose Gegenwart, Ewigkeit
\item Zufälligkeit in der Zeit; Nebeneinander im Raum
\item Fortgang, Entwicklung, Selbstkonstruktion
\item Reflexion und Anschauung
\item Aufbau in Stufen; Rand
\item Bestimmung durch Selbstbezug; Fixpunkt
\item Durchgang, Reihenfolge
\item Transzendentalphilosophie, Naturphilosophie
\item Gerade und Kreis; Umlauf
\item Kreisstrom, Magnetfeld, Pole
\item Sich-selbst-Anschauen; Ich = Ich
\item Sollen, Streben
\end{itemize}

\subsection*{4. Personen}
\begin{itemize}
\item Georg Wilhelm Friedrich Hegel
\item Johann Gottlieb Fichte
\item Friedrich Wilhelm Joseph Schelling
\item Immanuel Kant
\end{itemize}

\subsection*{5. Quellen}
\begin{itemize}
\item Hegel, G.\,W.\,F.: Differenz des Fichteschen und Schellingschen Systems der Philosophie (1801). In: Werke, Band~2: Jenaer Schriften 1801--1807. Frankfurt am Main: Suhrkamp 1986, S.~9--138. Zitiert als W~2.
\item Fichte, J.\,G.: Grundlage der gesamten Wissenschaftslehre; System der Sittenlehre. Nach Hegels Darstellung und Anmerkungen, W~2, S.~52~ff.
\end{itemize}

\subsection*{Dokumentdaten}
\begin{itemize}
\item Titel: Zeit und Ewigkeit der Rekursion. Der Fortgang der Selbstkonstruktion der Identität zur Totalität in Hegels Differenzschrift (1801)
\item Datei: zeit\_und\_ewigkeit\_der\_rekursion\_v1.tex
\item Version: v1
\item Datum: 2026-09-25
\item Herkunft: nimmt die in der Reflexion vom 2026-09-25 offen gelassene Frage nach Rekursion und Zeit auf; setzt die Essays \glqq Subjekt-Objekt und die Identität der Identität und der Nichtidentität\grqq{} (v2) und \glqq Zwei Rekursionen\grqq{} (v1) voraus
\item Rechnung: endliche Stufenbäume, selbstbezüglicher Baum, Fixpunktgleichung $s = 1 + 2s$; Python mit sympy, Ergebnisse in Abschnitt 3
\item Grundlage: Hegel, Werke Band~2 (Suhrkamp 1986), Textdatei der Vorlage; Seitenzahlen nach eingedruckter Paginierung
\item Status: Erstfassung; Kanon nach Pauschale, solange die Qualität hält
\end{itemize}

\end{document}
